% ============================================================
% content/solutions.tex —— 2025 湖南卷物理 · 参考答案与详细解析
% 解析类型：source_embedded（转录自官方 Word 版内嵌解析，独立复核）
% ============================================================

\begin{answer}
\gksolhead{1}{单项选择题}
\gkanswer{答案 B}
原子核衰变时释放能量，根据质能方程，总质量会减少，新核总质量小于原核质量，故 A 错误；半衰期定义为大量放射性原子核半数发生衰变所需的时间，题干中对数量的描述也是“大量”，符合定义，故 B 正确；半衰期由原子核内部结构决定，与温度无关，故 C 错误；半衰期只由核内部自身的因素决定，与原子所处的化学状态和外部条件没有关系，故 D 错误。
\end{answer}

\begin{answer}
\gksolhead{2}{单项选择题}
\gkanswer{答案 C}
由题意可知，物块沿斜面向上做匀减速直线运动，设初速度为 $v_{0}$，加速度大小为 $a$，斜面倾角为 $\theta$。水平方向上物块做匀减速直线运动，初速度为 $v_{0x}=v_{0}\cos\theta$，加速度大小为 $a_{x}=a\cos\theta$，则有
\[
  v_{x}^{2}-v_{0x}^{2}=-2a_{x}x,
\]
整理可得 $v_{x}=\sqrt{(v_{0}\cos\theta)^{2}-2a\cos\theta\cdot x}$，可知 $v_{x}$--$x$ 图像应为类似抛物线的一部分，故 A、B 错误；物块在竖直方向上做匀减速直线运动，速度为 $v_{0y}=v_{0}\sin\theta$，加速度大小为 $a_{y}=a\sin\theta$，则有
\[
  v_{y}^{2}-v_{0y}^{2}=-2a_{y}y,
\]
整理可得 $v_{y}=\sqrt{(v_{0}\sin\theta)^{2}-2a\sin\theta\cdot y}$，可知 $v_{y}$--$y$ 图像应为类似抛物线的一部分，故 C 正确，D 错误。

\begin{gktip}
【一题多解】物块沿光滑斜面上滑时，其水平方向和竖直方向的分运动均为匀变速直线运动，在匀变速直线运动中速度的二次方与位移成线性关系，故 $v_{x}$--$x$ 图像、$v_{y}$--$y$ 图像一定为曲线，通过排除法选 C。
\end{gktip}
\end{answer}

\begin{answer}
\gksolhead{3}{单项选择题}
\gkanswer{答案 D}
根据题意，画出光路图，如图甲所示，$O$ 为该介质横截面的圆心，由几何关系可知，折射角为 $45^{\circ}$，由折射定律有
\[
  n=\frac{\sin\theta}{\sin45^{\circ}}=\sqrt{2}\sin\theta>1,
\]
且 $\sin\theta<1$，则有 $n<\sqrt{2}$，$\sin\theta>\dfrac{\sqrt{2}}{2}$，即 $\theta>45^{\circ}$，故 A、B 错误；令 $\sin C=\dfrac{1}{n}$，根据折射光直接由 $B$ 点出射可知 $\sin C>\dfrac{\sqrt{2}}{2}$，即全反射临界角 $C>45^{\circ}$；增大入射角，光路图如图乙所示，由几何关系可知，光在 $\overset{\frown}{BC}$ 上的入射角小于 $45^{\circ}$，则该单色光在 $\overset{\frown}{BC}$ 上不可能发生全反射，故 C 错误；减小入射角，光路图如图丙所示，由几何关系可知，光在 $\overset{\frown}{AB}$ 上的入射角大于 $45^{\circ}$，可能大于临界角，则该单色光在 $\overset{\frown}{AB}$ 上可能发生全反射，故 D 正确。

\gkfig{figures/s3-optics}
\end{answer}

\begin{answer}
\gksolhead{4}{单项选择题}
\gkanswer{答案 A}
根据题意，卫星在小行星同步轨道和表面附近轨道运行时轨道半径分别为 $R+h$、$R$，设卫星的质量为 $m$，由开普勒第三定律可知
\[
  \frac{(R+h)^{3}}{T_{0}^{2}}=\frac{R^{3}}{T_{1}^{2}},
\]
解得 $R=\dfrac{T_{1}^{\frac{2}{3}}}{T_{0}^{\frac{2}{3}}-T_{1}^{\frac{2}{3}}}\,h$；卫星绕小行星表面附近做匀速圆周运动，由万有引力提供向心力可知
\[
  \frac{GMm}{R^{2}}=m\frac{4\pi^{2}}{T_{1}^{2}}R,
\]
解得 $M=\dfrac{4\pi^{2}R^{3}}{GT_{1}^{2}}$。对应结果可得 $a$ 为 $T_{1}$，$b$ 为 $T_{0}$，$c$ 为 $T_{1}$，故 A 正确。
\end{answer}

\begin{answer}
\gksolhead{5}{单项选择题}
\gkanswer{答案 C}
小球 $A$ 静止时，对小球 $A$ 受力分析，如图所示，由平行四边形定则及几何关系可知，轻绳上拉力为 $F_{T}=mg$，小球 $A$ 与小球 $B$ 间的库仑力 $F_{\text{库}}=2mg\cos30^{\circ}=\sqrt{3}\,mg$，故 A、B 错误；若将轻绳剪断，则剪断瞬间小球 $A$ 受到轻绳的拉力消失，其他两力保持不变，根据力的平衡可知，此时小球 $A$ 所受的合力大小为 $mg$，则加速度大小为 $g$，故 C 正确；若将轻绳剪断，则剪断瞬间小球 $B$ 受到的库仑力、重力不变，小球 $B$ 仍然处在静止状态，则轻杆对小球 $B$ 的作用力不变，故 D 错误。

\gkfig{figures/s5-forces}
\end{answer}

\begin{answer}
\gksolhead{6}{单项选择题}
\gkanswer{答案 B}
设变压器原、副线圈的电压分别为 $U_{1}$、$U_{2}$，电流分别为 $I_{1}$、$I_{2}$，电压关系 $\dfrac{U_{1}}{U_{2}}=k$，电流关系 $\dfrac{I_{1}}{I_{2}}=\dfrac{1}{k}$。在原线圈中 $U_{1}=U-I_{1}R_{\text{原}}$，在副线圈中 $U_{2}=I_{2}R$，故灯泡两端电压
\[
  U_{2}=\frac{U_{1}}{k}=\frac{U-I_{1}R_{\text{原}}}{k}
      =\frac{U}{k}-\frac{I_{2}R_{\text{原}}}{k^{2}}
      =\frac{U}{k}-\frac{U_{2}R_{\text{原}}}{k^{2}R},
\]
整理后有 $\dfrac{k^{2}R+R_{\text{原}}}{k^{2}R}U_{2}=\dfrac{U}{k}$，即灯泡两端电压 $U_{2}=\dfrac{kRU}{k^{2}R+R_{\text{原}}}$，故流过灯泡的电流为 $I_{2}=\dfrac{U_{2}}{R}=\dfrac{kU}{k^{2}R+R_{\text{原}}}$，灯泡的电功率
\[
  P=U_{2}I_{2}=\frac{k^{2}U^{2}R}{(k^{2}R+R_{\text{原}})^{2}}.
\]
由于不知道 $k$ 的具体值，故无法判断原线圈所在回路的阻值 $R_{\text{原}}$ 为何值时功率最大，故 A 错误；当开关 $S$ 与 $a$ 相连时，$R_{\text{原}}=3R$，故灯泡两端电压 $U_{2}=\dfrac{kU}{k^{2}+3}$，故 B 正确；当开关 $S$ 与 $b$ 相连时，$R_{\text{原}}=2R$，故流过灯泡的电流为 $I_{2}=\dfrac{kU}{(k^{2}+2)R}$，故 C 错误；当开关 $S$ 与 $c$ 相连时，$R_{\text{原}}=R$，故灯泡的电功率 $P=\dfrac{k^{2}U^{2}}{(k^{2}+1)^{2}R}$，故 D 错误。

\begin{gktip}
【一题多解】设理想变压器原、副线圈的匝数分别为 $n_{1}$、$n_{2}$，原、副线圈电压分别为 $U_{1}$、$U_{2}$，副线圈负载电阻为 $R$，可以将变压器与负载看为一个整体，等效为一个新的电阻
\[
  R_{\text{等}}=\frac{U_{1}}{I_{1}}
      =\frac{\frac{n_{1}}{n_{2}}U_{2}}{\frac{n_{2}}{n_{1}}I_{2}}
      =\left(\frac{n_{1}}{n_{2}}\right)^{2}R=k^{2}R,
\]
其中 $k=\dfrac{n_{1}}{n_{2}}$（原、副线圈匝数比），等效后电路如图所示。当开关 $S$ 与 $a$ 相连时，$R_{\text{原}}=3R$，等效电阻两端电压 $U_{1}=\dfrac{R_{\text{等}}}{R_{\text{等}}+3R}U=\dfrac{k^{2}}{k^{2}+3}U$，故灯泡两端电压 $U_{2}=\dfrac{U_{1}}{k}=\dfrac{kU}{k^{2}+3}$，故 B 正确；当开关 $S$ 与 $b$ 相连时，$R_{\text{原}}=2R$，流过等效电阻的电流 $I_{1}=\dfrac{U}{R_{\text{等}}+2R}=\dfrac{U}{(k^{2}+2)R}$，故流过灯泡的电流为 $I_{2}=kI_{1}=\dfrac{kU}{(k^{2}+2)R}$，C 错误；当开关 $S$ 与 $c$ 相连时，$R_{\text{原}}=R$，灯泡的电功率与等效电阻的电功率相同，流过等效电阻的电流 $I_{1}=\dfrac{U}{R_{\text{等}}+R}$，故功率 $P=I_{1}^{2}R_{\text{等}}=\dfrac{k^{2}U^{2}}{(k^{2}+1)^{2}R}$，故 D 错误；由闭合电路欧姆定律知识可知，当外电路电阻与内阻相等时功率最大，即 $R_{\text{原}}$ 和 $R_{\text{等}}$ 最接近时功率最大，由于 $k$ 未知，故无法判断开关 $S$ 与哪个触点接触时 $R_{\text{原}}$ 和 $R_{\text{等}}$ 最接近，即无法判断何时灯泡功率最大，故 A 错误。

\gkfig{figures/s6-equiv}
\end{gktip}
\end{answer}

\begin{answer}
\gksolhead{7}{多项选择题}
\gkanswer{答案 AD}
两横波的波长均为
\[
  \lambda=\frac{v}{f}=\SI{4}{\metre},
\]
故 A 正确；来自 $A$ 处、$B$ 处的两列波传到 $C$ 处所需时间分别为 $t_{1}=\dfrac{3}{10}\,\si{\second}=\SI{0.3}{\second}$，$t_{2}=\dfrac{\sqrt{3^{2}+4^{2}}}{10}\,\si{\second}=\SI{0.5}{\second}$，$T=\dfrac{1}{f}=\SI{0.4}{\second}$，故 $t=\SI{0.4}{\second}$ 时，$A$ 处波已传到 $C$ 处且 $C$ 处质点已经振动了 $\dfrac{1}{4}T$，故 $C$ 处质点处于正向或负向最大位移处，加速度最大，速度为 0，故 B、C 错误；分析可知 $t=\SI{0.6}{\second}$ 时两列波都已传播到 $C$ 处，$C$ 处质点到两波源的距离差为
\[
  \Delta x=5\,\si{\metre}-3\,\si{\metre}=2\,\si{\metre}=\frac{1}{2}\lambda,
\]
故 $C$ 处为振动减弱点，由于两列波振幅相同，故 $C$ 处位移始终为 0，速度为 0，故 D 正确。
\end{answer}

\begin{answer}
\gksolhead{8}{多项选择题}
\gkanswer{答案 AD}
根据题意可知 $O$ 点、$A$ 点和 $B$ 点的电势分别为
\[
  \varphi_{O}=\frac{E_{p}}{q},\qquad
  \varphi_{A}=-\frac{E_{p}}{q},\qquad
  \varphi_{B}=\frac{E_{p}}{2q},
\]
则 $OA$ 中点的电势为 $\varphi_{M}=\dfrac{\varphi_{O}+\varphi_{A}}{2}=0$，故 A 正确；如图所示，设 $N$ 点为 $AB$ 的三等分点，同理可知 $N$ 点电势为 0，连接 $MN$，则 $MN$ 为一条等势线，过 $A$ 点作 $MN$ 的垂线，可知电场线沿该垂线方向指向右下方，由几何关系可知 $\angle NMA=45^{\circ}$，故电场的方向与 $x$ 轴正方向成 $45^{\circ}$ 角，故 B 错误；电场强度的大小为
\[
  E=\frac{0-\left(-\dfrac{E_{p}}{q}\right)}{\dfrac{1}{2}d\cdot\cos45^{\circ}}
    =\frac{2\sqrt{2}\,E_{p}}{qd},
\]
故 C 错误，D 正确。

\gkfig{figures/s8-geometry}
\end{answer}

\begin{answer}
\gksolhead{9}{多项选择题}
\gkanswer{答案 AC}
根据右手定则可知金属杆沿 $x$ 轴正方向运动过程中，金属杆中电流沿 $y$ 轴负方向，故 A 正确；若金属杆可以在沿 $x$ 轴正方向的力 $F$ 作用下做匀速直线运动，则有
\[
  F=F_{\text{安}}=BIL,\qquad
  I=\frac{BLv_{0}}{R}=\frac{BLv_{0}}{Lr_{0}}=\frac{Bv_{0}}{r_{0}},
\]
可得 $F=\dfrac{B^{2}Lv_{0}}{r_{0}}$，由于金属杆运动过程中接入导轨中的长度 $L$ 在变化，故 $F$ 为变力，故 B 错误；取一微小时间 $\Delta t$，设此时金属杆接入导轨中的长度为 $L'$，根据动量定理有
\[
  -BI'L'\Delta t=-BL'q'=m\Delta v,
\]
同时有 $q=\dfrac{\dfrac{\Delta\Phi}{\Delta t}}{L'r_{0}}\cdot\Delta t=\dfrac{\Delta\Phi}{L'r_{0}}=\dfrac{B\cdot\Delta S}{L'r_{0}}$，联立得 $-\dfrac{B^{2}\Delta S}{r_{0}}=m\Delta v$，对从开始到金属杆停止运动时整个过程累积可得 $-\dfrac{B^{2}S}{r_{0}}=0-mv_{0}$，解得此时金属杆与导轨围成的面积
\[
  S=\frac{mv_{0}r_{0}}{B^{2}},
\]
故 C 正确；若金属杆的初速度减半，由上述分析可知当金属杆停止运动时金属杆与导轨围成的面积 $S'=\dfrac{1}{2}S$，根据抛物线的图像规律可知此时金属杆停止运动时经过的距离大于原来的一半，故 D 错误。
\end{answer}

\begin{answer}
\gksolhead{10}{多项选择题}
\gkanswer{答案 BD}
爆炸后，$A$、$B$ 组成的系统动量守恒，有 $3mv_{1}=mv_{2}$；$B$ 与 $C$ 碰撞过程动量守恒，有 $mv_{2}=6mv$，联立解得 $v=0.5v_{1}$。爆炸后瞬间 $A$ 的动能 $E_{kA}=\dfrac{1}{2}\cdot3m\cdot v_{1}^{2}$，$D$ 的初动能 $E_{kD}=\dfrac{1}{2}\cdot6m\cdot(0.5v_{1})^{2}$，两者不相等，故 A 错误；$D$ 水平滑动过程中摩擦力做功为
\[
  W_{f}=-\mu\cdot6mg\cdot s_{1}=-\mu\cdot6mg\cdot\frac{h}{\mu}=-6mgh,
\]
$D$ 做平抛运动过程中重力做的功为 $W_{G}=6mgh$，故 $D$ 从开始运动到落地瞬间合外力做功为 0，根据动能定理可知 $D$ 的初动能与其落地时的动能相等，故 B 正确；$D$ 物块平抛过程有 $h=\dfrac{1}{2}gt^{2}$，$s_{2}=v_{0}t$，联立解得 $v_{0}=s_{2}\sqrt{\dfrac{g}{2h}}$；$D$ 水平滑动过程中根据动能定理有
\[
  -6mgh=\frac{1}{2}\cdot6m\cdot v_{0}^{2}-\frac{1}{2}\cdot6mv^{2},
\]
化简得 $v^{2}=\dfrac{s_{2}^{2}g}{2h}+2gh$。弹药释放的能量完全转化为 $A$ 和 $B$ 的动能，则爆炸过程的能量为
\[
  E=\frac{1}{2}\cdot3m\cdot v_{1}^{2}+\frac{1}{2}\cdot m\cdot v_{2}^{2}
    =24mv^{2}
    =24m\left(\frac{s_{2}^{2}g}{2h}+2gh\right)
    =48mgh\left(1+\frac{s_{2}^{2}}{4h^{2}}\right),
\]
故 C 错误，D 正确。
\end{answer}

\begin{answer}
\gksolhead{11}{实验题}
\gkanswer{答案 (1) 2.206 mm（2.205--2.207 均可）；(3) 0.010 m/s}
\gkanswer{(4) $k=\dfrac{(\rho-\rho_{0})g\pi D^{2}}{6v}$；(5) 减小}
(1) 根据图 1 可知小球直径
\[
  D=2\,\si{\milli\metre}+20.6\times0.01\,\si{\milli\metre}=\SI{2.206}{\milli\metre}.
\]

(3) 由图 2 可知 $A$、$E$ 两点间的距离为 $x=(7.02-5.00)\times10^{-2}\,\si{\metre}=2.02\times10^{-2}\,\si{\metre}$，小球从 $A$ 运动到 $E$ 的时间为 $t=4t_{0}=4\times\SI{0.5}{\second}=\SI{2}{\second}$，所以速度为
\[
  v=\frac{x}{t}=\frac{2.02\times10^{-2}\,\si{\metre}}{2\,\si{\second}}\approx\SI{0.010}{\metre\per\second}.
\]

(4) 小球匀速运动，根据受力平衡有 $\rho gV=\rho_{0}gV+f$，由体积公式可知，小球的体积 $V=\dfrac{4}{3}\pi R^{3}=\dfrac{4}{3}\pi\left(\dfrac{D}{2}\right)^{3}$，整理可得
\[
  k=\frac{(\rho-\rho_{0})gV}{Dv}=\frac{(\rho-\rho_{0})g\pi D^{2}}{6v}.
\]

(5) 根据 (4) 可知匀速运动的速度与小球直径的二次方成正比，所以换成直径更小的同种材质小球，匀速运动时的速度将减小。
\end{answer}

\begin{answer}
\gksolhead{12}{实验题}
\gkanswer{答案 (1) 8000 Ω；(2) 见解析图；(3) $\rho=\dfrac{Rbc}{a}$}
\gkanswer{(5) $R_{2}$；(6) 大于}
(1) 由题可知多用电表开关旋转到“$\times$1 k”挡，故根据题图 1 可得读数为 $\SI{8000}{\ohm}$。

(2) 由上述知长方体导电水泥块样品的电阻 $R_{x}>\sqrt{R_{A}R_{V}}$，故采用电流表内接法；实验中要求滑动变阻器采用分压接法，故连接实物图如图所示。

\gkphoto[0.6\textwidth]{photos/image25-wiring.png}{伏安法接线图（原图保留，非矢量例外）}

(3) 电阻定律 $R=\dfrac{\rho L}{S}$，化简得 $\rho=\dfrac{RS}{L}$，根据题图 2 可知，测量的水泥块长度为 $a$，横截面积为 $bc$，代入解得
\[
  \rho=\frac{Rbc}{a}.
\]

(5) 根据题图 3 可知压力越大电阻率越小，即电阻越小；回路中电流增大，滑动变阻器 $R_{2}$ 电压增大，水泥块 $R_{1}$ 电压减小，而报警器在两端电压大于或等于 \SI{3}{\volt} 时启动，故应将报警器并联在 $R_{2}$ 两端。

(6) 电源电动势 $E$ 减小，要使报警器启动，即 $R_{2}$ 两端电压仍为 \SI{3}{\volt}，根据串联分压有
\[
  U_{2}=\frac{R_{2}}{R_{1}+R_{2}}E=\frac{1}{\dfrac{R_{1}}{R_{2}}+1}E,
\]
可知 $E$ 减小需要 $R_{1}$ 更小，又因为 $F$ 越大 $R_{1}$ 越小，可知 $F_{1}$ 大于 $F_{0}$。
\end{answer}

\begin{answer}
\gksolhead{13}{解答题}
\gkanswer{答案 (1) $g=\dfrac{p_{0}(L_{2}-L_{1})}{L_{1}\rho h}$}
\gkanswer{(2) $g=\SI{9.5}{\metre\per\second\squared}$}
(1) 竖直放置时空气柱的压强为 $p_{1}=p_{0}+\rho gh$；水平放置时空气柱的压强 $p_{2}=p_{0}$。由题可知，空气柱发生等温变化，根据玻意耳定律，有
\[
  p_{1}L_{1}S=p_{2}L_{2}S,
\]
解得
\[
  g=\frac{p_{0}(L_{2}-L_{1})}{L_{1}\rho h}.
\]

(2) 由题可知，空气柱发生等容变化，根据查理定律，有
\[
  \frac{p_{1}}{T_{1}}=\frac{p_{2}}{T_{2}},
\]
代入数据可得
\[
  g=\SI{9.5}{\metre\per\second\squared}.
\]
\end{answer}

\begin{answer}
\gksolhead{14}{解答题}
\gkanswer{答案 (1) $q=\dfrac{mv_{0}^{2}}{E_{0}}$}
\gkanswer{(2) $B=\dfrac{2E_{0}}{dv_{0}}$}
\gkanswer{(3) $x_{m}=\dfrac{(2+\sqrt{3})d}{2}$}
(1) 粒子在电容器中做类平抛运动，其中粒子在水平方向做匀速直线运动，有
\[
  \sqrt{3}\,d=v_{0}t,
\]
在竖直方向做匀变速直线运动，有
\[
  \frac{d}{2}=\frac{0+v_{y}}{2}t,\qquad
  v_{y}=at=\frac{qU}{md}t,
\]
由闭合回路欧姆定律可得
\[
  U=\frac{r_{0}}{r_{0}+2r_{0}}E_{0},
\]
联立解得 $v_{y}=\dfrac{\sqrt{3}}{3}v_{0}$，$q=\dfrac{mv_{0}^{2}}{E_{0}}$。

(2) 粒子进入磁场时的速度与竖直方向夹角的正弦值为 $\tan\theta=\dfrac{v_{x}}{v_{y}}$，即 $\theta=60^{\circ}$，$v=\dfrac{v_{0}}{\sin60^{\circ}}=\dfrac{2\sqrt{3}}{3}v_{0}$。粒子在进入磁场后做匀速圆周运动，即
\[
  qvB=m\frac{v^{2}}{R},
\]
其中由几何关系可得
\[
  R=\frac{\dfrac{d}{2}}{\cos30^{\circ}}=\frac{\sqrt{3}\,d}{3},
\]
联立解得
\[
  B=\frac{2E_{0}}{dv_{0}}.
\]

(3) 对粒子取一个竖直向上的速度使得该速度对应产生的洛伦兹力与水平向右的静电力平衡，则有
\[
  qv_{y1}B=qE,
\]
解得 $v_{y1}=\dfrac{2\sqrt{3}}{3}v_{0}$。此时粒子做圆周运动的合速度在竖直方向的分速度大小为
\[
  v_{y2}=v_{y1}'+v_{y}=\sqrt{3}\,v_{0},\qquad v_{y1}'=v_{y1}.
\]
此时合速度与竖直方向夹角的正切值为 $\tan\alpha=\dfrac{v_{0}}{\sqrt{3}\,v_{0}}$，合速度大小为 $v'=\sqrt{(\sqrt{3}\,v_{0})^{2}+v_{0}^{2}}$，粒子做圆周运动的半径 $r=\dfrac{mv'}{Bq}$，故最远距离为
\[
  x_{m}=r+r\cos\alpha=\frac{(2+\sqrt{3})d}{2}.
\]
\end{answer}

\begin{answer}
\gksolhead{15}{解答题}
\gkanswer{答案 (1) $F=4mg$；(2) $v=\sqrt{\dfrac{37gL}{10}}$}
\gkanswer{(3) $v=\sqrt{\dfrac{9kgL}{10(k+1)}+\dfrac{14}{5}gL}$（$k\ge1$），$v_{\min}=\sqrt{\dfrac{13}{4}gL}$}
(1) 对机器人由 $B$ 点运动到最低点的过程，由动能定理可得
\[
  \frac{1}{2}mv_{0}^{2}-\frac{1}{2}mv^{2}=mgL,
\]
对于机器人在最低点的运动状态，由向心力公式可得
\[
  F-mg=m\frac{v_{0}^{2}}{L},
\]
联立可得 $F=4mg$。

(2) 轻绳运动到左上方与水平方向夹角为 $37^{\circ}$ 时，由能量守恒定律可得
\[
  \frac{1}{2}mv^{2}=mgL\sin37^{\circ}+\frac{1}{2}mv_{2}'^{2},
\]
水平方向位移
\[
  x=L\cos37^{\circ}+L=v_{2}'\sin37^{\circ}\cdot t,
\]
竖直方向（取向上为正）位移
\[
  y=1.2L-L\sin37^{\circ}=v_{2}'\cos37^{\circ}\cdot t-\frac{1}{2}gt^{2},
\]
联立可得
\[
  v=\sqrt{\frac{37gL}{10}}.
\]

(3) 当机器人运动到滑杆左上方且与水平方向夹角为 $37^{\circ}$ 时，记此时机器人所在的点为 $C$ 点。在水平方向对机器人运动到 $C$ 点的过程，由动量守恒定律可得 $mx_{1}=Mx_{2}$，对该过程中机器人和杆的运动情况分析可得 $x_{1}+x_{2}=L+L\cos37^{\circ}$，解得
\[
  x_{1}=\frac{1.8kL}{1+k},\qquad x_{2}=\frac{1.8L}{1+k},
\]
$x_{1}$ 方向向左，$x_{2}$ 方向向右。令 $C$ 点到 $A$ 点的水平方向上的距离为 $x_{CA}$，对运动情况进行分析，可得
\[
  x_{CA}=L\cos37^{\circ}+(L-x_{2})=\frac{1.8kL}{1+k},
\]
同理，令 $C$ 点到 $A$ 点的竖直方向上的距离为 $y_{CA}$，则
\[
  y_{CA}=1.2L-L\cdot\sin37^{\circ}=0.6L.
\]
设此时机器人在水平方向的速度为 $v_{x}$，在竖直方向的速度为 $v_{y}$，杆的速度为 $v_{M}$，根据动量守恒定律可得 $mv_{x}=Mv_{M}$，即 $v_{x}=kv_{M}$。对于机器人由 $B$ 点运动到 $C$ 点的过程，由动能定理可得
\[
  \frac{1}{2}mv^{2}=\frac{1}{2}m(v_{x}^{2}+v_{y}^{2})
                  +\frac{1}{2}km\left(\frac{v_{x}}{k}\right)^{2}
                  +mgL\cdot\sin37^{\circ},
\]
对该式化简可得
\[
  v^{2}=\frac{k+1}{k}v_{x}^{2}+v_{y}^{2}+1.2gL.
\]
分析机器人由 $C$ 点运动到 $A$ 点的过程，可得
\[
  v_{x}\cdot t=\frac{1.8kL}{1+k}\ \text{（水平方向）},\qquad
  v_{y}\cdot t-\frac{1}{2}gt^{2}=0.6L\ \text{（竖直方向）}.
\]
当以杆为参考系时，小球依然在做圆周运动，此时小球相对于杆的水平速度 $v_{x}'=v_{x}+v_{M}$，由此可得 $v_{y}=\dfrac{v_{x}'}{\tan37^{\circ}}$，解得
\[
  v_{y}=\frac{4}{3}\cdot\frac{k+1}{k}\cdot v_{x}.
\]
联立 $v_{y}$ 与机器人从 $C$ 点到 $A$ 点竖直方向运动公式，可得
\[
  \frac{4}{3}\cdot\frac{k+1}{k}\cdot\frac{1.8kL}{1+k}-\frac{1}{2}gt^{2}=0.6L,
\]
联立解得 $t=\sqrt{\dfrac{3.6L}{g}}$。将 $t$ 代入机器人由 $C$ 点到 $A$ 点竖直方向运动的公式，可得
\[
  v_{y}\cdot\sqrt{\frac{3.6L}{g}}-\frac{1}{2}g\cdot\frac{3.6}{g}=0.6L,
\]
解得 $v_{y}=\sqrt{1.6gL}$。因为 $v_{y}=\dfrac{4}{3}\cdot\dfrac{k+1}{k}\cdot v_{x}$，可得 $v_{x}=\dfrac{3k}{4(k+1)}v_{y}$，即
\[
  v_{x}=\frac{3k}{4(k+1)}\cdot\sqrt{1.6gL}.
\]
将 $v_{x}$、$v_{y}$ 代入前面动能定理公式化简后得到的公式，可得
\[
  v^{2}=\frac{k+1}{k}\cdot\frac{9k^{2}}{16(k+1)^{2}}\cdot1.6gL
       +1.6gL+1.2gL,
\]
解得
\[
  v=\sqrt{\frac{9kgL}{10(k+1)}+\frac{14}{5}gL}\qquad(k\ge1).
\]
当 $k=1$ 时，$v$ 值最小，为
\[
  v_{\min}=\sqrt{\frac{13gL}{4}}.
\]
\end{answer}
